% ---------------------------------------------------------------
\section{Discussion}\label{sec:discussion}
% ---------------------------------------------------------------

\subsection{What an auditor may conclude}

\begin{figure}[t]
\centering
\resizebox{\columnwidth}{!}{% Fig.: what an auditor may conclude (Discussion), column width
\begin{tikzpicture}[font=\sffamily\scriptsize, >={Stealth[length=1.6mm]},
  q/.style={draw=black!55, fill=black!4, rounded corners=2pt, text width=24mm, align=center,
            inner sep=2pt, minimum height=13mm,
            execute at begin node={\hyphenpenalty=10000\exhyphenpenalty=10000}},
  a/.style={draw=black!30, fill=white, rounded corners=2pt, text width=47mm, align=flush left,
            inner sep=2.5pt, minimum height=13mm,
            execute at begin node={\hyphenpenalty=10000\exhyphenpenalty=10000}},
  fl/.style={->, black!65, line width=0.6pt}]
\newcommand{\Y}{\textcolor{okgreen}{\textbf{yes}}}
\newcommand{\N}{\textcolor{advorange}{\textbf{no}}}
\def\dy{1.65}
\node[q] (q1) at (0,0)      {\textbf{1} Is the trainer trusted?};
\node[q] (q2) at (0,-\dy)   {\textbf{2} Does the verifier fix the seed after the data commitment?};
\node[q] (q3) at (0,-2*\dy) {\textbf{3} Is there an authenticated route from sources to training?};
\node[q] (q4) at (0,-3*\dy) {\textbf{4} Is the model free of poison?};
\node[a, right=3mm of q1] (a1) {\Y: attested capture binds the model to its record; no proof of training needed.\\
  \N: binding (exact replay or a proof of training) is required. \hfill\textcolor{black!55}{Thm.~\ref{thm:bind}}};
\node[a, right=3mm of q2] (a2) {\Y: no prover influence, if replay is exact on a pinned stack.\\
  \N: up to $\log_2 T$ chosen bits for $T$ runs. \hfill\textcolor{black!55}{Thms.~\ref{thm:noinf},~\ref{thm:grind}; G2, G5}};
\node[a, right=3mm of q3] (a3) {\Y, with binding: lineage holds via source witnesses or attested capture and transform; stage types must be declared. \N: no lineage claim. \hfill\textcolor{black!55}{Thm.~\ref{thm:auth}, Cor.~\ref{cor:launder}}};
\node[a, right=3mm of q4, draw=advorange, fill=advorange!7] (a4) {Never certified. At best: every poisoned item used was missed by the detector, worth its recall, which collapsed once the backdoor worked. \hfill\textcolor{black!55}{Thm.~\ref{thm:content}; G3, G5}};
\foreach \i [count=\j from 2] in {1,2,3} \draw[fl] (q\i) -- (q\j);
\foreach \i in {1,...,4} \draw[fl] (q\i) -- (a\i);
\end{tikzpicture}
}
\caption{Four questions an auditor should answer before accepting a claim about a released
model, and what the theorems and experiments say for each answer.}
\label{fig:auditor}
\end{figure}

Table~\ref{tab:minimal} can be read as a checklist, and Fig.~\ref{fig:auditor} condenses it
into four questions. Before accepting a claim about a
released model, an auditor should establish three things:
\begin{enumerate}
\item Who is trusted: which profile applies.
\item Which mechanisms were actually deployed and verified.
\item Which goal the claim corresponds to.
\end{enumerate}
Two examples show how this works.

A cloud provider's attested data pipeline without binding (\textsf{AttCap}, \textsf{AttTrans}
under the full adversary) supports no claim about the released model. It does support the
record-level claim $G_{\mathrm{compl}}$.

A signed-manifest scheme with a trusted training infrastructure (the trans-only or cap-only
profiles) does support lineage for filter-only pipelines, but not the absence of poison.

\subsection{Answers to five design questions}

\paragraph{Binding}
Binding is required whenever the trainer is not trusted, and no amount of witnessing or
attestation of earlier stages substitutes for it (Theorem~\ref{thm:bind}). If binding is
used, its randomness must be fixed after commitment (Theorems~\ref{thm:noinf}
and~\ref{thm:grind}). Exact re-execution is a valid binding only where training is
bit-reproducible. Section~\ref{sec:eval} measures when it is.

\paragraph{Witnessing}
Source receipts must be content-addressed (split view, Section~\ref{sec:eval}). Source
witnessing and input witnessing answer different questions, lineage and completeness. Each
fails under collusion, so witnesses should be organisationally independent of the operator
they observe.

\paragraph{Synthetic data}
Generative stages need an attested transform even on the witness route, and stage types
must be declared and authenticated (Corollary~\ref{cor:launder}). Lineage of synthetic data
is recursive: the generator is itself a model whose provenance matters. Our model treats it
as an authorised but possibly poisoned source (\textsf{AdvGen}).

\paragraph{Blockchain versus transparency logs}
For every goal we study, a witnessed transparency log is sufficient as the record store.
A ledger adds inclusion liveness (Proposition~\ref{prop:ledger}). We find no support for
the claim that a blockchain makes training records complete or truthful. Those properties
are decided upstream of the store.

\paragraph{Poisoning}
Provenance cannot certify that a model is poison-free (Theorem~\ref{thm:content}). It can
make a detector's verdict \emph{transfer} to the released model, which is what binding and
completeness contribute. The resulting claim is only as strong as the detector's recall.
In our experiments, recall was weakest exactly where the poison worked. The prototype's detector flagged at most a quarter of the planted items. On DistilBERT, recall fell from 68--83\% at the lower, ineffective poison rates to 25\% at the 2\% rate where the backdoor took hold (Section~\ref{sec:eval}). Any
certification scheme that reports provenance without detector recall overstates its
assurance.

\subsection{Open problems}

\begin{enumerate}
\item \emph{Binding at pre-training scale.} Zero-knowledge binding now reaches LoRA
  fine-tuning~\cite{liao2026verilora}, but not full pre-training. Exact re-execution
  doubles the cost and requires determinism. Spot-check proofs of learning admit spoofing
  under tolerance~\cite{fang2023pol}.
\item \emph{Tolerance as a grinding channel.} Theorem~\ref{thm:grind} bounds influence by the
  number of accepted trajectories. For tolerance-based verifiers that number is not
  characterised. Optimum vicinity~\cite{tan2025zkpot} closes the channel for convex
  objectives; for non-convex LLM training it remains open.
\item \emph{Attested accelerators.} The attestation routes assume that the relevant stages run
  in TEEs. Atlas notes that GPUs lack equivalent protection~\cite{spoczynski2025atlas}.
\item \emph{Detector guarantees.} $G_{\mathrm{content}}$ is parameterised by recall. Certified
  recall bounds for poisoning detectors at LLM scale are lacking.
\item \emph{Confidential verifiability ($P_7$).} All our mechanisms reveal the record to the
  auditor. Combining them with confidentiality of training data is open.
\item \emph{Composition across model generations.} Fine-tuning chains and synthetic data make
  provenance recursive. The model covers one training run.
\end{enumerate}

\subsection{Limitations}

\paragraph{Model}
The model has five abstractions:
\begin{itemize}
\item Items are atomic and content-identified.
\item Each mechanism is all-or-nothing.
\item A single training run is modelled.
\item Detectors are abstracted to an examined set and a miss set, with no budget or base
  rate.
\item Correlated compromise is captured \emph{across} trust anchors (any subset of
  capabilities), but not \emph{within} a mechanism class, such as one corrupt witness among
  several.
\end{itemize}
Table~\ref{tab:scope} lists which training practices fall inside the abstraction. The
model's results are exact for the model only. Their relevance to real systems rests on the
validation in Section~\ref{sec:eval}.

\paragraph{Prototype}
The prototype trains a small bag-of-words network on 2,000 reviews. It simulates trusted
execution with key isolation rather than real enclaves, so it validates the logic of
attestation-based checks, not any hardware. It compares three detectors on one small model. Poisoning at GPU scale was
measured on one model (DistilBERT), one dirty-label trigger and three seeds per condition,
with one detector in G3 and three in G5, so it is reported descriptively. G5 ran on a second,
smaller GPU because the first was unavailable; its qualitative results match G3, but its runs
did not replay G3's bit for bit. Bit-exact replay therefore holds on a fixed stack only, and
we changed GPU and library versions together, so we cannot say which change breaks it. The
GPU experiments cover fine-tuning, not pre-training.
Their seed counts (160 and 48) limit the largest $T$ we could test to 128 and 32.

\paragraph{Literature}
The literature search used one database, title-and-abstract screening was performed by an
AI assistant, and full-text assessment is pending (Section~\ref{sec:lit}). The case-study
mapping relies partly on abstracts.

\paragraph{Proofs}
The proofs are by hand. They are backed by two independent exhaustive decision procedures
for the six profiles, and by a third, vectorised checker over all 128 adversary subsets. They
are not mechanised in a proof assistant. Proposition~\ref{prop:ledger} and the composition
remark of Section~\ref{sec:scope} are argued, not proved.
